\section{引言}
约 2 bit/parameter 的压缩需要向量、加性码本、格码或 trellis 等结构化表示~\cite{gptvq,aqlm,quipsharp,qtip}。但部署 checkpoint 并非单一整数张量：block-scaled VQ 同时保留共享码本、组尺度和离散索引。固定码本后，尺度是组级连续坐标，assignment 是向量级离散坐标。PV-Tuning 已将量化表示中的连续/离散变量形式化为交替坐标优化~\cite{pvtuning}；本文不声称首次提出坐标交替，而追问更窄的问题：在固定共享码本、只允许当前码字附近领域切换时，两种坐标能否在同一软松弛中直接组合？

同锚点实验先揭示功能非均匀性：尺度适应把六任务平均提高 3.72 个百分点，但 PPL 退化 4.28\%；当前码字 8-way 邻域 assignment 仅接受 0.378\% 的硬切换，提高 1.48 个百分点，PPL 退化 2.27\%。随后，最直接的同步联合训练只得到 69.74\%/11.4098，并被尺度端点在准确率与 PPL 上同时支配。更关键的是，joint 过程的尺度配合零切换时甚至低于未修改锚点。这将失败定位为软码字 mixture 上的补偿与硬部署之间的失配，而非泛化的优化不稳定。

本文贡献如下：第一，在同一 2-bit checkpoint、监督、层和码率下，证明组尺度与局部 assignment 具有不同任务--语言建模效应；第二，完整验证并否决 naive simultaneous joint，并用零切换硬状态定位补偿--投影失配；第三，导出部署态交接原则。由该原则产生的分阶段修复仍待正式验证，本文不提前宣称完成的联合方法。
